Um die Transformationsverhalten der gegebenen bilinearen Kovarianten unter Lorentztransformationen zu beweisen, verwenden wir die Transformationseigenschaften der Dirac-Spinoren und der Gamma-Matrizen.

Zunächst erinnern wir uns daran, dass ein Dirac-Spinor $\psi$ unter einer Lorentztransformation $\Lambda$ wie folgt transformiert:

\[
\psi \rightarrow S(\Lambda) \psi
\]

wobei $S(\Lambda)$ die Spinor-Repräsentation der Lorentztransformation ist. Die adjungierte Spinor $\bar{\psi}$ transformiert entsprechend:

\[
\bar{\psi} \rightarrow \bar{\psi} S^{-1}(\Lambda)
\]

Nun betrachten wir die einzelnen bilinearen Formen:

1. **Pseudoskalar $\bar{\psi} \gamma^5 \psi$:**

Die Transformation dieser Größe ergibt sich aus:

\[
\bar{\psi} \gamma^5 \psi \rightarrow \bar{\psi} S^{-1}(\Lambda) \gamma^5 S(\Lambda) \psi
\]

Unter Verwendung der gegebenen Transformationseigenschaft von $\gamma^5$, nämlich $S^{-1}(\Lambda) \gamma^5 S(\Lambda) = \operatorname{det}(\Lambda) \gamma^5$, erhalten wir:

\[
\bar{\psi} \gamma^5 \psi \rightarrow \operatorname{det}(\Lambda) \bar{\psi} \gamma^5 \psi
\]

Dies zeigt, dass $\bar{\psi} \gamma^5 \psi$ als Pseudoskalar transformiert.

2. **Axialer Vektor $\bar{\psi} \gamma^\mu \gamma^5 \psi$:**

Für diesen Ausdruck haben wir:

\[
\bar{\psi} \gamma^\mu \gamma^5 \psi \rightarrow \bar{\psi} S^{-1}(\Lambda) \gamma^\mu \gamma^5 S(\Lambda) \psi
\]

Da $\gamma^\mu$ als Vektor transformiert, gilt:

\[
S^{-1}(\Lambda) \gamma^\mu S(\Lambda) = \Lambda_\nu^\mu \gamma^\nu
\]

Kombiniert mit der Transformationseigenschaft von $\gamma^5$, ergibt sich:

\[
\bar{\psi} \gamma^\mu \gamma^5 \psi \rightarrow \operatorname{det}(\Lambda) \Lambda_\nu^\mu \bar{\psi} \gamma^\nu \gamma^5 \psi
\]

Dies zeigt, dass $\bar{\psi} \gamma^\mu \gamma^5 \psi$ als axialer Vektor transformiert.

3. **Antisymmetrischer Tensor $\bar{\psi} \sigma^{\mu \nu} \psi$:**

Für den antisymmetrischen Tensor $\sigma^{\mu \nu} = \frac{\mathrm{i}}{2} [\gamma^\mu, \gamma^\nu]$ ergibt sich die Transformation:

\[
\bar{\psi} \sigma^{\mu \nu} \psi \rightarrow \bar{\psi} S^{-1}(\Lambda) \sigma^{\mu \nu} S(\Lambda) \psi
\]

Da $\sigma^{\mu \nu}$ aus den $\gamma$-Matrizen besteht, die als Tensoren transformieren, folgt:

\[
S^{-1}(\Lambda) \sigma^{\mu \nu} S(\Lambda) = \Lambda_\alpha^\mu \Lambda_\beta^\nu \sigma^{\alpha \beta}
\]

Daraus ergibt sich:

\[
\bar{\psi} \sigma^{\mu \nu} \psi \rightarrow \Lambda_\alpha^\mu \Lambda_\beta^\nu \bar{\psi} \sigma^{\alpha \beta} \psi
\]

Dies zeigt, dass $\bar{\psi} \sigma^{\mu \nu} \psi$ als antisymmetrischer Tensor transformiert. 

Zusammenfassend haben wir gezeigt, dass die bilinearen Formen die angegebenen Transformationseigenschaften unter Lorentztransformationen besitzen.